\section{Background and terminology}
\label{sec:background}
\subsection{Synthetic data for vision}
\draftnote{Define synthetic images and labels, rendering/simulation, learned
generation and augmentation. Introduce classification, detection and
segmentation and their annotation needs. Define the intended real-world
deployment population before discussing quality.}

\subsection{Quality dimensions}
For this review, realism denotes perceptual or physical plausibility;
similarity denotes agreement with a specified reference distribution;
coverage concerns representation of relevant classes and conditions;
and bias concerns systematic imbalances and differences in errors across
relevant groups or conditions. These are working definitions to refine
against the selected literature. Downstream \gls{utility} concerns model
performance on the intended task after training with synthetic data.

\textcite{alaaHowFaithfulYour2022} distinguish fidelity, diversity and
generalization using $\alpha$-Precision, $\beta$-Recall and Authenticity.
This is a methodological starting point; vision-specific suitability
remains a question for the synthesis.
\draftnote{Distinguish image, annotation, dataset and trained-model
properties, and a requirement from its metric and assessment resources.}
